\phantomsection
\addcontentsline{toc}{section}{Source Code Scripts}
\section*{Source Code Scripts}

\subsection*{A.1 C\# MQTT Middleware Gateway (Batching Logic)}
The following code snippet demonstrates the internal concurrent queue and batch-flush mechanism responsible for mitigating extreme network load spikes by pooling PostgreSQL insertion requests.

\begin{lstlisting}[language=Csh, caption=Middleware Batched Ingestion Logic (C\#)]
// Program.cs - MqttToMiddlewareClient
using MQTTnet;
using Supabase;
using System.Collections.Concurrent;

class Program
{
    private static ConcurrentQueue<SensorData> _messageQueue = new();
    private static Client _supabaseClient = null!;
    private static bool _isRunning = true;

    // Background Flush Task (The "Shock Absorber")
    private static async Task PeriodicallyFlushToSupabase()
    {
        const int BATCH_SIZE = 100;
        const int FLUSH_INTERVAL_MS = 2000;

        while (_isRunning || !_messageQueue.IsEmpty)
        {
            if (_messageQueue.Count >= BATCH_SIZE || (_messageQueue.Count > 0 && _isRunning))
            {
                var listToInsert = new List<SensorData>();
                while (listToInsert.Count < BATCH_SIZE && _messageQueue.TryDequeue(out var data))
                {
                    listToInsert.Add(data);
                }

                if (listToInsert.Any())
                {
                    try
                    {
                        Console.WriteLine($"Flushing {listToInsert.Count} records to Supabase...");
                        await _supabaseClient.From<SensorData>().Insert(listToInsert);
                    }
                    catch (Exception ex)
                    {
                        Console.WriteLine($"Flush Error: {ex.Message}");
                    }
                }
            }

            await Task.Delay(FLUSH_INTERVAL_MS);
        }
    }
}
\end{lstlisting}

\newpage
\subsection*{A.2 Python Database Benchmark Orchestrator}
This script utilizes the Supabase REST API via PostgREST to objectively retrieve server-side execution planning metrics without introducing local network request latency.

\begin{lstlisting}[language=Python, caption=PostgREST Benchmark Orchestrator (Python)]
import requests
import re
import csv

SUPABASE_URL = "https://[YOUR_INSTANCE].supabase.co"
ANON_TOKEN = "[REDACTED_API_TOKEN]"
ITERATIONS = 100

def get_internal_execution_time(table_type):
    headers = {
        "apikey": ANON_TOKEN,
        "Authorization": f"Bearer {ANON_TOKEN}",
        "Content-Type": "application/json"
    }
    
    url = f"{SUPABASE_URL}/rest/v1/rpc/get_exact_benchmark_plan"
    
    try:
        response = requests.post(url, headers=headers, json={"tbl_type": table_type})
        if response.status_code != 200:
            return None
            
        plan_text = response.json()
        match = re.search(r"Execution Time: ([\d\.]+) ms", plan_text)
        if match:
            return float(match.group(1))
        return None
        
    except Exception as e:
        return None

def main():
    stats = []
    for i in range(1, ITERATIONS + 1):
        t_std = get_internal_execution_time("standard")
        t_hyp = get_internal_execution_time("hypertable")
        
        if t_std and t_hyp:
            stats.append({"run": i, "standard_ms": t_std, "hypertable_ms": t_hyp})

    with open('supabase_explain_results.csv', 'w', newline='') as f:
        writer = csv.DictWriter(f, fieldnames=["run", "standard_ms", "hypertable_ms"])
        writer.writeheader()
        writer.writerows(stats)

if __name__ == "__main__":
    main()
\end{lstlisting}

\subsection*{A.3 PostgreSQL RPC: Exact Benchmark Execution}
This server-side script guarantees identical workload formulation between standard table sets and the chunked hypertable environment.

\begin{lstlisting}[language=SQL, caption=PostgreSQL EXPLAIN ANALYZE Loop]
CREATE OR REPLACE FUNCTION public.get_exact_benchmark_plan(tbl_type text)
RETURNS text AS $$
DECLARE
  line text;
  full_plan text := '';
  sql_q text;
BEGIN
  IF tbl_type = 'standard' THEN
    sql_q := 'SELECT device_id, AVG(moisture) 
              FROM public.standard_telemetry_table 
              WHERE created_at > NOW() - INTERVAL ''7 days'' 
              GROUP BY device_id';
  ELSE
    sql_q := 'SELECT device_id, AVG(moisture) 
              FROM public.telemetry_hypertable 
              WHERE time > NOW() - INTERVAL ''7 days'' 
              GROUP BY device_id';
  END IF;
  
  FOR line IN EXECUTE 'EXPLAIN (ANALYZE) ' || sql_q LOOP
    full_plan := full_plan || line || chr(10);
  END LOOP;
  RETURN full_plan;
END;
$$ LANGUAGE plpgsql SECURITY DEFINER;
\end{lstlisting}
